% ===================================================================
%  TABEL Nr. 12 — Jaam. --- Raudteed. --- Laevad.
%  Traduction 2026 d'apres texto/io, controlee sur texto/fr et
%  texto/fi. Consigne : voir texto/et/10-tabel-01.tex.
%
%  LA CHOSE LA PLUS NEUVE DE CETTE PAGE PORTE LE MOT LE PLUS VIEUX,
%  ET CE MOT VIENT DE L'EMPIRE MONGOL.
%
%  Le chemin de fer est « raudtee », la route-de-fer, compose fait a
%  la maison sur « raud » le fer et « tee » la route. Rien la que de
%  banal : la chose etait nouvelle, on l'a decrite.
%
%  MAIS LA GARE EST « JAAM », ET « JAAM » EST LE RUSSE « ЯМ » — le
%  relais de poste a chevaux que les Mongols avaient organise en
%  travers de l'Asie et dont la Russie a herite le reseau avec le
%  nom. Ce mot a traverse un continent avant d'arriver sur un quai
%  balte, et il y designe aujourd'hui un batiment de brique du XIXe
%  siecle.
%
%  POURQUOI L'UN A-T-IL ETE COMPOSE ET L'AUTRE HERITE ? Parce que la
%  voie ferree ne remplacait rien — personne n'avait jamais eu de
%  route de fer — tandis que LA GARE REMPLACAIT QUELQUE CHOSE, et
%  precisement le relais de poste : le meme batiment au bord de la
%  meme route, ou l'on attendait, ou l'on changeait de vehicule, ou
%  l'on faisait peser ses malles.
%
%  ET C'EST UN TROISIEME MECANISME, A COTE DE L'EMPRUNT ET DE LA
%  COMPOSITION : L'HERITAGE. UNE CHOSE NEUVE QUI PREND LA PLACE D'UNE
%  ANCIENNE PREND SON NOM, QUELLE QUE SOIT SON ORIGINE. Le tableau 12
%  romanche disait : qui fonde nomme, et la Viafier Retica en etait la
%  preuve. L'estonien ajoute le cas ou personne ne nomme, parce que le
%  nom etait deja la et qu'il a suffi de le laisser sur place.
%
%  Le reste de la page se partage comme on l'attend maintenant :
%  « rong » le train, « vedur » la locomotive de « vedama » tirer,
%  « pilet », « perroon », « konduktor » et « lokomobiil » visiblement
%  etrangers, et « raudtee », « veetorn », « kaubarong » composes.
% ===================================================================

%%K t12-apar-1 apar
\VUsaut{4.0mm}
\VUtitre{15.0pt}{60}{TABEL Nr. 12}
\VUsaut{2.4mm}
\VUpk{11.4pt}{40}{Jaam. --- Raudteed. --- Laevad.}
\VUsaut{3.4mm}
\textsc{Märkus.} --- \textit{Igas riigis kasutatavad sõiduplaanid on erinevad, me
kasutame seetõttu näitena} \VUgras{prantsuse tabeli} \textit{kellaaegu.}

%%K t12-01-1 p
1. --- Ma reisisin äsja läbi Saksamaa. Enne väljasõitu ostsin ma \VUgras{sõiduplaani}
\textsuperscript{(15)}, et teada rongide väljumisaegu. Olles kindlaks teinud, et kõige
sobivam rong väljub Pariisist kell üheksa õhtul ja saabub Strasbourgi kell üks
kolmteist minutit, valmistasin ma oma reisi ette, ja järgmisel päeval lasin ma end
jaama viia.

%%K t12-02-1 p
2. --- Kui ma astusin suurde väljumissaali, vana maa\VUgras{preestri}
\textsuperscript{(2)} järel, kes kandis käes oma \VUgras{reisikotti}
\textsuperscript{(3)}, oli palju \VUgras{reisijaid} \textsuperscript{(1)} juba kohal.

%%K t12-02-2 p
Kuna ma tahtsin saata \VUgras{telegrammi} \textsuperscript{(5)}, lasin ma
\VUgras{kontrolöril} \textsuperscript{(6)} endale näidata \VUgras{telegraafikontorit}
\textsuperscript{(4)}.

%%K t12-03-1 p
3. --- Enne \VUgras{ooteruumi} \textsuperscript{(7)} minekut läksin ma end varustama
\VUgras{piletiga} \textsuperscript{(9)}, milles oli ka tagasisõit.

%%K t12-04-1 p
4. --- Ma ei oodanud kaua \VUgras{luugi} \textsuperscript{(8)} juures, sest seal oli vähe
inimesi. \VUgras{Sõdur} \textsuperscript{(10)}, kes oli puhkusele minemas, oli nii lahke
ja loovutas mulle oma järjekorra, nii et mul oli küllalt aega küsida mõnda teadet
\VUgras{jaamaülemalt} \textsuperscript{(13)}, kes vestles järelevalve \VUgras{komissari}
\textsuperscript{(12)} ja valves oleva \VUgras{sandarmiga} \textsuperscript{(11)}.

%%K t12-tit-1 sub
\VUpk{10.2pt}{40}{Pagasi registreerimine.}
\VUsaut{3.0mm}

%%K t12-05-1 p
5. --- Olles ostnud ajalehe \VUgras{raamatukioskist} \textsuperscript{(14)}, lasin ma
kaaluda oma \VUgras{pagasi} \textsuperscript{(17)}, mille \VUgras{kandja}
\textsuperscript{(16)} oli äsja toonud \VUgras{pakivankriga} \textsuperscript{(19)};
\VUgras{ametnik} \textsuperscript{(22)}, kellele oli usaldatud \VUgras{kaal}
\textsuperscript{(21)}, teatas mulle, et mu kohver kaalub kolmkümmend viis kilogrammi;
see tegi viis kilogrammi \VUgras{ülekaalu}, mille eest ma pean maksma lisatasu
\VUgras{pagasiluugi} \textsuperscript{(23)} juures.

%%K t12-06-1 p
6. --- Kuna ma olin liiga vara, oli mul vaba aega vaadelda reisijate liikumist jaamas.
Ühed andsid ära või võtsid tagasi oma väikese pagasi \VUgras{pakihoiuruumis}
\textsuperscript{(25)}; teised uurisid \VUgras{sõiduplaane} \textsuperscript{(26)}.
Saabuvad reisijad kiirustasid \VUgras{väljapääsu} \textsuperscript{(32)} poole oma
\VUgras{kohvriga} \textsuperscript{(24)} käes. Mõned ootasid \VUgras{pagasi väljastuse}
\textsuperscript{(27)} juures ja esitasid oma \VUgras{kviitungi} \textsuperscript{(29)},
et tagasi saada oma kohvreid, \VUgras{kaste} \textsuperscript{(28)} või \VUgras{korve}
\textsuperscript{(30)}.

%%K t12-07-1 p
7. --- \VUgras{Aktsiisiametnikud} \textsuperscript{(33)} uurisid hoolikalt pakke, et
kindlaks teha, kas on midagi deklareerida.

%%K t12-08-1 p
8. --- Mõned reisijad läksid \VUgras{einelauda} \textsuperscript{(34)}, kus
\VUgras{kelner} \textsuperscript{(35)} nägi vaeva neid teenindades. Nad peavad
kiirustama, sest \VUgras{rong} \textsuperscript{(36)}, mis on kiirrong, ei seisa kaua
jaamas.

%%K t12-09-1 p
9. --- Siis läksin ma väljumisperroonile ja otsisin otsest \VUgras{vagunit}
\textsuperscript{(37)}, et ei oleks vaja teel vagunit vahetada. Ma astusin
koridorvagunisse, milles on \VUgras{kupee} \textsuperscript{(38)} suitsetajatele ja kupee,
mis on reserveeritud «üksi reisivatele daamidele».

%%K t12-tit-2 sub
\VUpk{10.2pt}{40}{Öine väljasõit.}
\VUsaut{3.0mm}

%%K t12-10-1 p
10. --- Olles üles astunud \VUgras{astmelauale} \textsuperscript{(40)}, sulgenud
\VUgras{ukse} \textsuperscript{(39)}, üles kerinud \VUgras{ruloo} \textsuperscript{(41)} ja
pannud oma väikese pagasi \VUgras{võrgule} \textsuperscript{(43)}, istusin ma
\VUgras{pingile} \textsuperscript{(42)}. Nähes \VUgras{lambisüütajat}
\textsuperscript{(44)} lampe süütamas, mõtlesin ma, et pean veetma öö vagunis, ja ma
kutsusin ametniku, kelle hooleks on \VUgras{peapatjade} \textsuperscript{(45)}
väljaüürimine, et üht paluda.

%%K t12-11-1 p
11. --- Konduktor tuli pileteid kontrollima. Ma küsisin temalt, miks me veel ei ole
välja sõitnud, sest kell on õige. Ta vastas mulle, et \VUgras{rongiülem}
\textsuperscript{(46)} on ametis jalgrataste panemisega \VUgras{pakivagunisse}
\textsuperscript{(47)}. Pealegi ei olnud veel kõiki jalasoojendajaid
\textsuperscript{(48)} vahetatud.

%%K t12-12-1 p
12. --- Kuid me pidime peagi väljuma, sest \VUgras{kütja} \textsuperscript{(50)} võttis
oma \VUgras{tendrilt} \textsuperscript{(49)} sütt, et ergutada \VUgras{veduri}
\textsuperscript{(54)} tuld, ja \VUgras{vedurijuht} \textsuperscript{(51)} ootas ainult
\VUgras{jaamaülema abi} \textsuperscript{(52)} signaali, kes oli \VUgras{perroonil}
\textsuperscript{(53)}.

%%K t12-13-1 p
13. --- Veduri \VUgras{katel} \textsuperscript{(56)} mürises tuimalt;
\VUgras{korsten} \textsuperscript{(55)} oksendas suitsujugasid, mis olid segunenud
\VUgras{auruga} \textsuperscript{(60)}. Kiledalt kõlas \VUgras{vile}
\textsuperscript{(59)} ja \VUgras{kolvid} \textsuperscript{(58)} hakkasid liikuma,
lükates raske masina \VUgras{rattaid} \textsuperscript{(57)}.

%%K t12-tit-3 sub
\VUsaut{3.0mm}
\VUpk{10.2pt}{40}{Väljasõit!}
\VUsaut{3.0mm}

%%K t12-14-1 p
14. --- Rongi kiirus \VUgras{rööbastel} \textsuperscript{(64)} kasvas;
\VUgras{perroonid} \textsuperscript{(65)} kadusid; me möödusime \VUgras{käimlatest}
\textsuperscript{(66)} ja ka vagunitest, mis olid pandud teisele \VUgras{teele}
\textsuperscript{(63)}: \VUgras{magamisvagunid} \textsuperscript{(67)},
\VUgras{restoranivagun} \textsuperscript{(68)}, \VUgras{postivagun}
\textsuperscript{(69)}.

%%K t12-noto-1 noto
\VUnotes{143.2mm}{%
\textit{(*) Märkus.} --- \textit{Muidugi on reisikirjeldus tehtud sõjaeelse piiri
järgi.}%
}

%%K t12-15-1 p
15. --- Siis möödus meie silme eest \VUgras{kaubajaam} \textsuperscript{(71)}. Angaaride
all laadisid ametnikud \VUgras{kaupa} \textsuperscript{(72)}, mis saadeti aeglase rongiga.
\VUgras{Manöövrimehed} \textsuperscript{(76)} \VUgras{veetorni} \textsuperscript{(73)}
juures manööverdasid \VUgras{loomavaguneid} \textsuperscript{(75)} ja \VUgras{lahtiseid
vaguneid} \textsuperscript{(79)}, et moodustada \VUgras{kaubarongi}
\textsuperscript{(80)}.

%%K t12-16-1 p
16. --- Me ületasime viimase \VUgras{pöörangu} \textsuperscript{(78)}, mille juures
seisis pööranguvaht; me möödusime \VUgras{signaalkettast} \textsuperscript{(86)} ja jaam
kadus kaugusesse.

%%K t12-17-1 p
17. --- Peagi ületasime me \VUgras{raudsilla} \textsuperscript{(74)}, mis läheb üle
\VUgras{jõe} \textsuperscript{(81)}. Olles selle silla ületanud, sõitsime me mööda jõge,
millel võimsad \VUgras{puksiirid} \textsuperscript{(82)} vedasid raskeid \VUgras{praame}
\textsuperscript{(83)}. Me nägime veel \VUgras{reisilaeva} \textsuperscript{(84)}, kui see
oli randumas \VUgras{pontooni} \textsuperscript{(85)} juurde.

%%K t12-18-1 p
18. --- Olles suure kiirusega möödunud mitmest jaamast, läbisime me mõned
\VUgras{tunnelid} \textsuperscript{(87)} ja jätsime enda taha loendamatud tõkkepuuvahtide
\VUgras{majakesed} \textsuperscript{(88)}. Rongi õõtsumisest hällitatuna jäin ma sügavasse
unne.

%%K t12-tit-4 sub
\VUsaut{3.0mm}
\VUpk{10.2pt}{40}{Deutsch-Avricourti jaam.}
\VUsaut{3.0mm}

%%K t12-19-1 p
19. --- Mind äratas järsku ametniku hääl, kes hüüdis: «Deutsch-Avricourt!
\textsuperscript{(*)}. Kõik välja!» Toimus tollikontroll ja kõik piiri tüütud
formaalsused. Ma pidin seadma oma taskukella jaama \VUgras{suure kella}
\textsuperscript{(89)} järgi, mis oli viiskümmend viis minutit ees; vaevalt ärganuna
eristasin ma raskesti \VUgras{suurt osutit} \textsuperscript{(90)} \VUgras{väiksest}
\textsuperscript{(91)}; \VUgras{numbrilaud} \textsuperscript{(92)} ise oli hommikuse udu
tõttu täiesti segane.

%%K t12-20-1 p
20. --- Sellal kui tolliametnikud tegid oma kontrolli, lugesin ma haigutades
\VUgras{plakateid} \textsuperscript{(93)}, mis kaunistavad jaama seinu. Ma sõin aja
veetmiseks hommikust, uurisin Saksa \VUgras{võrgu} \textsuperscript{(94)} kaarti, et näha,
kui palju maad mind veel Strasbourgist lahutab, ja läksin tagasi oma vagunisse,
tugevdatuna lootusest peagi oma reisi siht saavutada.

\VUsaut{6.0mm}
\VUfilet{22mm}
